% ======================================================================================
% RESULTS -- R1..R5.  Budget ~3000 words across 14 main-text figures, ~215 words each.
%
% Section titles are deliberately NOT claims. An earlier draft called R5 "a representation
% problem"; the evidence supports "a representation-DEPENDENT failure", which is weaker and is
% what the titles below say. Do not strengthen a heading without a figure that carries it.
% ======================================================================================

\section{Results}
\label{sec:results}

% --------------------------------------------------------------------- R1
\subsection{Two scenarios, two forgetting phenotypes}
\label{sec:r1}

\wordcount{450}
% Fig.~\ref{fig:911} is the hinge of the whole report: it is the evidence for splitting the
% mechanism investigation by SCENARIO rather than by RULE. Read the columns first (they
% differ), then the rows (they barely do).
%
% Numbers, all ten seeds, matched competence, re-derived from the saved arrays:
%   final task 1     Class-IL  5.4    Domain-IL 38.5
%   crossover        Class-IL 65.0    Domain-IL 75.8
%   ceilings         Class-IL 93.6    Domain-IL 94.3   -- retention is read against these
%
% The STRUCTURAL half of the argument matters more than the numbers and must come first:
% Class-IL has 10 output units, five of which receive no positive target during task 2, so
% output suppression is AVAILABLE. Domain-IL has 5 shared units and suppression CANNOT occur.
% The scenarios do not differ in degree; they differ in which mechanisms are physically
% possible. Only the output layer differs between them -- same rule, same protocol, same seeds.
%
% Be honest about the half that does not split: the argmax-minus-probe gap is large in BOTH
% (+81.6 and +29.5), even though suppression is impossible in one. R5 returns to this.

\begin{figure*}[t]
  \fig{911_two_forgetting_phenotypes.png}
  \caption{\label{fig:911}\wordcount{130}}
\end{figure*}

% --------------------------------------------------------------------- R2
\subsection{The size of the learning-rule effect}
\label{sec:r2}

\wordcount{650}
% Order matters here. Fig.~\ref{fig:912} answers "is it real"; Fig.~\ref{fig:913} answers "is
% it large". Answering them in the other order makes the first look like special pleading.
%
% At the working configuration (H = 32, one hidden layer):
%   PC - backprop crossover   Class-IL  +1.42 +- 0.39   (9W-1L, p = 0.021)
%                             Domain-IL -0.72 +- 0.19   (0W-10L, p = 0.002)
% Small, real, and scenario-flipped. Task-1 phase lengths are comparable (719/707 vs 755/787
% steps), so this is NOT PC retaining by training less -- say so explicitly, it is the first
% objection a reader will raise.
%
% Depth widens the split in BOTH directions at once (+3.44 vs -3.49 at four layers), and width
% adds a reversal: Class-IL H = 4 is -4.14 +- 0.70, PC strongly WORSE at the narrowest width.
% That survives the dt correction (it was -6.19 under the confounded dt = 0.4).
%
% ⚠ The lr sweep must be reported with its censoring, not through it: at lr 0.16 the Class-IL
% crossover is undefined on 10/10 backprop seeds and 6/10 PC seeds. There is no measurement
% there. Fig.~\ref{fig:912}'s lower row carries this and the text must not average over it.
%
% Then the scale. Replay: +9.84 +- 0.77 and +3.43 +- 0.50 -- about seven times PC's effect
% where PC wins, and positive where PC is negative, flat across every lr, width and depth.
% The conclusion to state plainly: the problem is solvable, and changing the credit-assignment
% rule is not what solves it. Note also what replay costs -- it stores and re-presents task-1
% data, the exact resource a continual-learning rule is meant not to need. The comparison is
% of effect size under a fixed protocol, not of cost.

\begin{figure*}[t]
  \fig{912_pc_minus_backprop_sweeps.png}
  \caption{\label{fig:912}\wordcount{150}}
\end{figure*}

% figure*, not figure: panel (b) is saved at 7.06in (full \textwidth) because it carries three
% sweeps side by side, and it cannot be included in a 3.40in column without scaling -- which
% would shrink its 8pt labels. Panel (a) is column-width and is centred on its own row.
\begin{figure*}[t]
  \begin{subfigure}{\textwidth}
    \centering
    \fig{913_replay_sets_the_scale_a.png}
    \subcaption{\label{fig:913a}\wordcount{45}}
  \end{subfigure}

  \begin{subfigure}{\textwidth}
    \centering
    \fig{913_replay_sets_the_scale_b.png}
    \subcaption{\label{fig:913b}\wordcount{55}}
  \end{subfigure}
  \caption{\label{fig:913}\wordcount{50}}
\end{figure*}

% --------------------------------------------------------------------- R3
\subsection{Testing the mechanism prospective configuration is credited with}
\label{sec:r3}

\wordcount{500}
% A DIRECT test of the source paper's own claim (Song \& Bogacz Fig. 3b), and it belongs in
% Results whichever way it falls -- an earlier draft demoted it to an appendix negative, which
% was wrong: a negative result on the mechanism a paper credits is a result about that paper.
%
% Two quantities from the same instrumented run (script 803):
%   target alignment   cos(target - out_before, out_after - out_before), through training and
%                      against retention. Also in the interference mode (ref=), measured on a
%                      fixed task-1 batch DURING task 2 -- which is the quantity that matters
%                      for forgetting and which the source paper does not report.
%   displacement       D = ||x*(settled) - x(feedforward)||, non-zero at equilibrium BY
%                      CONSTRUCTION. Say so, or a non-zero D reads as evidence when it is a
%                      definition. The measurement is whether D predicts ||dW|| and retention.

\begin{figure}[t]
  \fig{915_target_alignment.png}
  \caption{\label{fig:915}\gap{915 not written; 803 has run}\wordcount{110}}
\end{figure}

\begin{figure}[t]
  \fig{916_displacement_to_retention.png}
  \caption{\label{fig:916}\gap{916 not written; 803 has run}\wordcount{110}}
\end{figure}

% --------------------------------------------------------------------- R4
\subsection{What repeated alternation reveals about the trajectory}
\label{sec:r4}

\wordcount{700}
% ⚠ THE PRE-REGISTERED PREDICTION LOOKS WRONG AND THE SECTION MUST SAY SO. The plan predicted
% Class-IL traces a closed LOOP under repeated alternation while Domain-IL SPIRALS IN. 804's
% block lengths do not support that as stated: Class-IL block lengths collapse hard (backprop
% seed 18: 530 -> 40 updates across ten blocks, ending task 1 at 46.3) against ~5% after a
% single switch, so Class-IL is converging too.
%
% But it is NOT uniform and the exceptions are the interesting part: Class-IL PC seed 13 hit
% the 5000-update cap twice and ended at task-1 0.2, where its neighbours end near 45.
% Domain-IL PC seeds 12 and 13 keep LONG, irregular task-2 blocks (up to 3870) while their
% task-1 blocks shrink.
%
% Hold the claim until Fig.~\ref{fig:917} draws the actual trajectory: block length and
% end-of-block accuracy are not the same thing as trajectory contraction. If the shape holds,
% the honest claim is that the scenarios differ in HOW FAST they find a joint solution, not in
% whether they can -- which is a more interesting result than the prediction was.
%
% Then 919: joint pre-training protects only PARTIALLY (Class-IL joint-start retains
% 6.9/25.6/12.5/20.3 against scratch 1.5/6.2/3.1/9.1, still collapsing from ~77%). So
% forgetting is not simply a failure to FIND a joint solution -- the solution can be handed to
% the network and it still leaves.

\begin{figure*}[t]
  \fig{917_alternation_geometry.png}
  \caption{\label{fig:917}\gap{917 not written; 804 Class-IL done, Domain-IL running}\wordcount{150}}
\end{figure*}

\begin{figure}[t]
  \fig{918_weight_space_pca.png}
  \caption{\label{fig:918}\gap{918 not written; needs 804}\wordcount{110}}
\end{figure}

\begin{figure}[t]
  \fig{919_joint_then_sequential.png}
  \caption{\label{fig:919}\gap{919 not written; 805 queued}\wordcount{110}}
\end{figure}

\begin{figure}[t]
  \fig{920_what_separates_by_scenario.png}
  \caption{\label{fig:920}\gap{920 not written; re-analysis only}\wordcount{110}}
\end{figure}

% --------------------------------------------------------------------- R5
\subsection{Evidence for an output-competition component}
\label{sec:r5}

\wordcount{700}
% Title deliberately weak. An earlier draft said "an output and calibration problem"; the
% evidence supports a COMPONENT, not the whole account, and the tension below is why.
%
% THE TENSION IS THE FINDING AND THIS SECTION MUST STAGE IT, NOT RESOLVE IT:
%   masking the absent classes recovers +17.5 points (script 120)
%   freezing the whole output layer recovers NOTHING, +0.36 +- 0.43 (script 130)
% Both spare the task-1 output weights during task 2, so naively they should agree.
%
% Script 806 separates them, and its smoke result landed on the SECOND pre-committed branch:
% control 3.8, freeze-all-W2 6.4, freeze task-1 columns 7.2, mask 65.0. Freezing exactly the
% weights masking spares does NOT reproduce masking -- and adding the task-1 biases moved it
% only 6.5 -> 7.2, so that is not the explanation either.
% ⚠ IF THIS SURVIVES TEN SEEDS, REWRITE THIS SECTION. Masking would then not work by sparing
% those weights, and the remaining candidate is that masking changes what W1 learns -- the
% trunk reshaped to serve the suppression objective -- which is not a readout account at all.
%
% The probe (Fig.~\ref{fig:923}) must draw its RANDOM-INIT FLOOR or it over-claims exactly as
% NCM did. Measured at H = 32: untrained probe 81.8%, after task 1 86.2%, against argmax
% 12.2 -> 91.8. Only ~4 points of the probe's range are attributable to training the trunk.
% NCM is excluded and the reason is stated: it has almost no dynamic range here.

\begin{figure}[t]
  \fig{921_retention_distribution.png}
  \caption{\label{fig:921}\gap{921 not written}\wordcount{110}}
\end{figure}

\begin{figure}[t]
  \fig{922_partial_column_freeze.png}
  \caption{\label{fig:922}\gap{922 not written; 806 queued}\wordcount{110}}
\end{figure}

\begin{figure*}[t]
  \fig{923_probe_vs_argmax.png}
  \caption{\label{fig:923}\gap{923 not written; 803 has run}\wordcount{150}}
\end{figure*}

\begin{figure}[t]
  \fig{924_weight_path_and_code_drift.png}
  \caption{\label{fig:924}\gap{924 not written; needs 803 and 804}\wordcount{110}}
\end{figure}

\begin{figure}[t]
  \fig{925_task_structure.png}
  \caption{\label{fig:925}\gap{925 not written}\wordcount{110}}
\end{figure}
